\documentclass[11pt,a4paper]{article}
\usepackage[utf8]{inputenc}
\usepackage[margin=1in]{geometry}
\usepackage{amsmath,amssymb,amsthm}
\usepackage{listings}
\usepackage{xcolor}
\usepackage{hyperref}

\newtheorem{theorem}{Theorem}
\newtheorem{lemma}[theorem]{Lemma}

\lstset{
  language=C++,
  basicstyle=\ttfamily\small,
  keywordstyle=\color{blue}\bfseries,
  commentstyle=\color{gray},
  stringstyle=\color{red},
  numbers=left,
  numberstyle=\tiny\color{gray},
  breaklines=true,
  frame=single,
  tabsize=4,
  showstringspaces=false
}

\title{IOI 2015 -- Sorting}
\author{}
\date{}

\begin{document}
\maketitle

\section{Problem Summary}

Given a permutation $S[0..n{-}1]$ and a sequence of $m$ grader swaps $(P_i, Q_i)$, in each round $i$ the grader first applies swap $(P_i, Q_i)$ to $S$, then you apply your own swap $(X_i, Y_i)$. Find the minimum number of rounds $m^*$ to sort $S$, and output your swaps.

\section{Solution}

\subsection{Binary Search on the Number of Rounds}

\begin{lemma}
The minimum number of swaps to sort a permutation equals $n$ minus the number of cycles in the permutation.
\end{lemma}

Let $S'$ be the permutation obtained by applying grader swaps $(P_0,Q_0), \ldots, (P_{m-1}, Q_{m-1})$ to $S$. Then $m$ rounds suffice if and only if the number of swaps needed to sort $S'$ is at most $m$, i.e., $n - \text{cycles}(S') \le m$.

This predicate is monotone in $m$ (applying more grader swaps can only be compensated by more of our swaps), so we binary search on $m$.

\subsection{Constructing the Swap Sequence}

Given the optimal $m^*$:
\begin{enumerate}
  \item Compute $S'$ by applying grader swaps $0, \ldots, m^*{-}1$ to $S$.
  \item Find the cycle decomposition of $S'$ and generate a list of $\le m^*$ swaps that sort $S'$ (pad with dummy swaps $(0,0)$).
  \item These swaps are expressed in terms of \emph{elements} (not positions). Simulate forward: at each round, after the grader's swap, find the current positions of the two target elements and swap those positions.
\end{enumerate}

\section{C++ Implementation}

\begin{lstlisting}
#include <bits/stdc++.h>
using namespace std;

int findSwapPairs(int n, int S[], int m,
                  int P[], int Q[], int X[], int Y[]) {
    auto check = [&](int rounds) -> bool {
        vector<int> T(S, S + n);
        for (int i = 0; i < rounds; i++) swap(T[P[i]], T[Q[i]]);
        vector<bool> vis(n, false);
        int cycles = 0;
        for (int i = 0; i < n; i++) {
            if (!vis[i]) {
                cycles++;
                for (int j = i; !vis[j]; j = T[j]) vis[j] = true;
            }
        }
        return n - cycles <= rounds;
    };

    int lo = 0, hi = m;
    while (lo < hi) {
        int mid = (lo + hi) / 2;
        if (check(mid)) hi = mid; else lo = mid + 1;
    }
    int ans = lo;

    // Build swap sequence
    vector<int> T(S, S + n);
    for (int i = 0; i < ans; i++) swap(T[P[i]], T[Q[i]]);

    vector<pair<int,int>> planned;
    vector<bool> vis(n, false);
    for (int i = 0; i < n; i++) {
        if (!vis[i] && T[i] != i) {
            vector<int> cyc;
            for (int j = i; !vis[j]; j = T[j]) { vis[j] = true; cyc.push_back(j); }
            for (int k = (int)cyc.size()-1; k >= 1; k--)
                planned.push_back({cyc[0], cyc[k]});
        }
    }
    while ((int)planned.size() < ans) planned.push_back({0, 0});

    // Simulate forward, translating element-swaps to position-swaps
    vector<int> pos(n), elem(n);
    for (int i = 0; i < n; i++) { elem[i] = S[i]; pos[S[i]] = i; }

    for (int i = 0; i < ans; i++) {
        // Grader swap
        int a = elem[P[i]], b = elem[Q[i]];
        swap(elem[P[i]], elem[Q[i]]);
        pos[a] = Q[i]; pos[b] = P[i];

        // Our swap (by element)
        int u = planned[i].first, v = planned[i].second;
        X[i] = pos[u]; Y[i] = pos[v];
        swap(elem[X[i]], elem[Y[i]]);
        pos[u] = Y[i]; pos[v] = X[i];
    }
    return ans;
}
\end{lstlisting}

\section{Complexity Analysis}

\begin{itemize}
  \item \textbf{Time}: $O((n + m) \log m)$ for binary search, $O(n + m)$ for construction.
  \item \textbf{Space}: $O(n + m)$.
\end{itemize}

\end{document}
